% Nouveau EGA V (cote 25): archive cote 25, site page 44 (PDF page 45).
% typescript. Draft compared: none.
% AI-transcribed by Claude Opus 5.5 (claude-opus-5-5, adaptive thinking, effort high), Message Batches API, batch ega25b (msgbatch_01A7n5YuXubUPNK1zsepi8RG), 2026-11-13; original AI transcription, subsequently checked against both supplied scan parts; see PAGES_CHECKED.tsv and CHECK_DECISIONS.tsv.
\scanpage{25/44}
% NOTE: The upper foliation is a legible final40 over unreadable earlier traces. The first morphism glyph remains mu/psi doubtful; the two later occurrences read mu.

\mspage{\hand{40}}

dans T tel que \hand{$\xi$} \var{\inserted{$\hand{\in} U(k)$}}{$\hand{\in} U(k)$}{insertion}{ega25_25_044-v0001} \var{\deleted{U}}{}{deletion}{ega25_25_044-v0002} implique qu'il existe exactement  un point \var{\xout{du}}{}{typewriter strike-out}{ega25_25_044-v0003} y de
$\underline{Y}^{\text{sing}}(k)$ au dessus de \hand{$\xi$}. Cela résultera de 8.\var{\over{\unsure{6}}{7}}{7}{writing over}{ega25_25_044-v0004}. c) qui implique plus
\var{\xout{généralement}}{}{typewriter strike-out}{ega25_25_044-v0005} précisément:

\uline{Corollaire} 8.\var{\over{\unsure{8}}{9}}{9}{writing over}{ega25_25_044-v0006}.  Supposons la condition (i) de 8.\var{\over{\unsure{7}}{8}}{8}{writing over}{ega25_25_044-v0007}. satisfaite, et
soit  \var{\xout{\unsure{T(k) un point tel que T soit lisse en} \illegible{1} U \unsure{Il en semble}}}{}{typewriter strike-out}{ega25_25_044-v0008}
U l'ouvert de T des points en lesquels T est lisse sur k. Alors \var{\inserted[typed]{$U\hand{\neq} 0$, \var{\deleted{\hand{et}}}{}{deletion}{ega25_25_044-v0010}}}{$U\hand{\neq} 0$, \var{\deleted{\hand{et}}}{}{deletion}{ega25_25_044-v0010}}{insertion}{ega25_25_044-v0009}
$\underline{Y}^{\text{sing}}|U \hand{\longrightarrow} U$ est une immersion ouverte, \var{\xout{(et même un isomorphisme}}{}{typewriter strike-out}{ega25_25_044-v0011} \var{\inserted{\unsure{à} fortiori $\underline{Y}^{\text{sing}}|U$ ne contient pas des points de $\underline{Y}^{\text{sup sing}}$.}}{\unsure{à} fortiori $\underline{Y}^{\text{sing}}|U$ ne contient pas des points de $\underline{Y}^{\text{sup sing}}$.}{insertion}{ega25_25_044-v0012}
\var{\xout{si X est propre sur k donc $\underline{Y}^{\text{sing}}$xxxx T surjectif)}}{}{typewriter strike-out}{ega25_25_044-v0013} Si X est propre sur
k donc $g:\underline{Y}^{\text{sing}}\hand{\longrightarrow} T$ surjectif, alors $\underline{Y}^{\text{sing}}|U\hand{\longrightarrow} U$ est un isomorphi
me, et U est le p\unsure{l|p}us grand ouvert de T ayant cette dernière propriété.

Tout d'abord, comme g est dominant et gén \var{\deleted{ramifié}}{}{deletion}{ega25_25_044-v0014} \var{\inserted{étale}}{étale}{insertion}{ega25_25_044-v0015}, il est
gén. étale, donc on peut \unsure{trouver|tppuver} un ouvert non vide V de T tel que
\var{\xout{\unsure{gén on peut supposer de plus g étale au(x)} \illegible{1} \unsure{que} \illegible{1} \unsure{ne contienne pas} $\underline{Y}$ \illegible{1}}}{}{typewriter strike-out}{ega25_25_044-v0016}
$\underline{Y}^{\text{sing}}|V \hand{\longrightarrow} V$ soit étale et surjectif\var{\inserted{\var{\margin{ce qui implique V lisse$/k$}}{\placedmargin{ce qui implique V lisse$/k$}}{marginal note placed}{ega25_25_044-v0018}}}{\var{\margin{ce qui implique V lisse$/k$}}{\placedmargin{ce qui implique V lisse$/k$}}{marginal note placed}{ega25_25_044-v0018}}{insertion}{ega25_25_044-v0017}. Si alors \hand{$\xi$}$\hand{\in} V(k)$ et si y est
un point de $\underline{Y}^{\text{sing}}(k)$ au dessus de \hand{$\xi$}, alors avec les notations de
8.\var{\over{\unsure{6}}{7}}{7}{writing over}{ega25_25_044-v0019}. c) \var{\xout{\unsure{Xlex}}}{}{typewriter strike-out}{ega25_25_044-v0020} l'espace H' n'est autre que l'espace tangent à T en \hand{$\xi$},
et comme on a noté à cet endroit, cela implique que le point x de X(k)
projection de y est déterminé comme le point orthogonal à H', donc est
uniquement déterminé, donc comme $\underline{Y}^{\text{sing}}\hand{\subset} X\text{x}$$\overset{\hand{\vee}}{P}$, y est uniquement déter-%
miné. Cela prouve déjà que g est birationnel (étant \var{\inserted[typed]{\var{\deleted{dominant}}{}{deletion}{ega25_25_044-v0022},}}{\var{\deleted{dominant}}{}{deletion}{ega25_25_044-v0022},}{insertion}{ega25_25_044-v0021} \var{\xout{\unsure{non ramifié}}}{}{typewriter strike-out}{ega25_25_044-v0023}
génériquement \var{\deleted{non ramffié}}{}{deletion}{ega25_25_044-v0024} \var{\inserted{étale}}{étale}{insertion}{ega25_25_044-v0025} et gén radiciel). \var{\xout{Pour Evidemment V H}}{}{typewriter strike-out}{ega25_25_044-v0026}
D'autre part, le morphisme \hand{$\unsure{\mu|\psi}$}(dont la définition en forme est évidente)
qui associe à tout \var{\over{x}{$\xi$}}{$\xi$}{writing over}{ega25_25_044-v0027}$\hand{\in} U(k)$ l'unique point $x=$ \hand{$\mu$}(\hand{$\xi$}) $\hand{\in} P$ orthogonal à
l'espace tangent à U en \hand{$\xi$}, coïncide sur V avec le composé
$V\hand{\rightarrow} \underline{Y}^{\text{sing}}|V \hand{\longrightarrow} X$, où la deuxième flèche est la projection; donc
\var{\inserted[typed]{posant}}{posant}{insertion}{ega25_25_044-v0028} \var{\xout{La restriction de}}{}{typewriter strike-out}{ega25_25_044-v0029} $h=($\hand{$\mu$}$,\mathrm{id}) : U\hand{\rightarrow} P\text{x}T$ , \var{\xout{\unsure{(il induit une immersion à V et est}}}{}{typewriter strike-out}{ega25_25_044-v0030}
\var{\xout{\unsure{montre que le composé j g}}}{}{typewriter strike-out}{ega25_25_044-v0031} $g_1 = g|g^{-1}(U)$\var{\xout{=}}{}{typewriter strike-out}{ega25_25_044-v0032}: $g^{-1}(U)\hand{\rightarrow} U$\var{\xout{\unsure{Ex}}}{}{typewriter strike-out}{ega25_25_044-v0033} , le composé

\var{\margin{\hand{44}}}{\placedmargin{\hand{44}}}{marginal note placed}{ega25_25_044-v0034}
